\section{Symmetric return maps and quartic pencils}
\label{sec:symmetric-models}

The four distances admit coordinates in which their common squareclass
and the final rational square are separate equations.  We first give the
return map, including its normalization, and then use the two rulings of
the radius quadric to obtain pencils of plane quartics.

Put
\[
 \Phi(u)=u^4-6u^2+1,
 \qquad
 \begin{aligned}
 N_0&=1-h-k-hk,&N_1&=1-h+k+hk,\\
 N_2&=1+h-k+hk,&N_3&=1+h+k-hk.
 \end{aligned}
\]

\begin{theorem}[Symmetric return]
\label{thm:sym-return}
Let $h,k,W\in\Q$ satisfy $W^2=\Phi(h)\Phi(k)$ and set
\[
 A=(1+h^2)(1+k^2),\qquad T=A+W.
\]
If $T\ne0$, the point
\begin{equation}
 x=\frac12-\frac{2k(1-h^2)}T,
 \qquad
 y=\frac12-\frac{2h(1-k^2)}T
 \label{eq:sym-return}
\end{equation}
satisfies the four distance identities $q_j=N_j^2/T$ in the corner
order $00,10,01,11$.

Conversely, let $(x,y)$ be a real point different from the corners,
let $r_j>0$ be its four distances, and put $S=r_0+r_1+r_2+r_3$.
Then
\begin{equation}
 h=\frac{r_2+r_3-r_0-r_1}{S},\qquad
 k=\frac{r_1+r_3-r_0-r_2}{S},\qquad
 T=\frac{16}{S^2},\qquad W=T-(1+h^2)(1+k^2)
 \label{eq:sym-forward}
\end{equation}
satisfy $W^2=\Phi(h)\Phi(k)$, $N_j=4r_j/S$, and
\eqref{eq:sym-return} returns the same point.  If
$r_j\in\sqrt d\,\Q_{>0}$ for one positive squarefree $d$, all four
parameters are rational and $[T]=[d]$.

In particular, the strict rational-distance problem is equivalent to
the existence of rational $h,k,W$ for which
\begin{equation}
 \begin{gathered}
 W^2=\Phi(h)\Phi(k),\qquad
 T=(1+h^2)(1+k^2)+W\in\Q^{\times2},\\
 4|k(1-h^2)|<T,\qquad 4|h(1-k^2)|<T.
 \end{gathered}
 \label{eq:sym-complete-system}
\end{equation}
\end{theorem}

\begin{proof}
The identity
\begin{equation}
 A^2-\Phi(h)\Phi(k)
 =8\bigl(k^2(1-h^2)^2+h^2(1-k^2)^2\bigr)
 \label{eq:sym-discriminant}
\end{equation}
gives the first distance equation on substituting
$xT=T/2-2k(1-h^2)$ and $yT=T/2-2h(1-k^2)$.
Indeed, twice its residual is
$T^2-2AT+A^2-\Phi(h)\Phi(k)=W^2-\Phi(h)\Phi(k)$.
The other three follow from
\[
 N_1^2-N_0^2=4k(1-h^2),\quad
 N_2^2-N_0^2=4h(1-k^2),\quad
 N_0^2+N_3^2=N_1^2+N_2^2.
\]

For the converse, $r_0^2+r_3^2=r_1^2+r_2^2$.  Expanding the
definitions in \eqref{eq:sym-forward} using this relation gives
$S^2N_j=4r_jS$.  For example,
\[
 \frac{S^2N_0}{2}
 =r_0^2+r_1^2+r_2^2-r_3^2+2r_0(r_1+r_2+r_3)
 =2r_0S.
\]
Thus $q_j=N_j^2/T$.  The differences
$q_1-q_0=1-2x$ and $q_2-q_0=1-2y$ give the return map.
Substitution in $q_0=x^2+y^2$ and
\eqref{eq:sym-discriminant} gives $W^2=\Phi(h)\Phi(k)$.
If $r_j=\sqrt d\,\rho_j$ with positive rational $\rho_j$, then
$S=\sqrt d\,s$ for $s\in\Q_{>0}$, so all the asserted
rationalities and $[T]=[d]$ follow.

Finally, when $T=t_0^2\ne0$, the return gives positive rational
distances $|N_j/t_0|$ whenever $0<x,y<1$; an interior point has no
zero distance.  The last two inequalities in
\eqref{eq:sym-complete-system} are exactly $|1-2x|<1$ and
$|1-2y|<1$.  The forward construction proves the reverse implication.
\end{proof}

The excluded case $T=0$ is explicit.  By
\eqref{eq:sym-discriminant}, it requires
$k(1-h^2)=h(1-k^2)=0$, hence
$(h,k)=(0,0)$ or $(h,k)=(\pm1,\pm1)$, with $W=-A$.
Also, $\Phi$ has no rational root.  These observations account for
the exceptional denominators in the return map.

\subsection{The homogeneous squareclass}

For a positive squarefree $n$, consider the intersection
\[
 D_n\subset\PP^3_{[a:b:c:e]}:\qquad
 a^2=e^2+c^2,\qquad nb^2=e^2-c^2.
\]

\begin{proposition}[Normalization of the return]
\label{prop:sym-normalization}
For $n>1$, every rational point of $D_n$ has $a+e\ne0$.
The functions
\[
 h=\frac{c}{a+e},\qquad \beta=\frac{2b}{a+e}
\]
identify its rational points with $n\beta^2=\Phi(h)$.
For two such points the returned squareclass is the class of
\begin{equation}
 T(p,q)=\frac{4(aa'+nbb')}{(a+e)(a'+e')}.
 \label{eq:sym-projective-T}
\end{equation}
This expression is invariant under independent rescaling of both
projective representatives.

The curve has no rational point if $n$ is even.  If $n$ is odd and
$(a,b,c,e)$ is a primitive integral representative with $a>0$, then
\[
 a+e=
 \begin{cases}
 2s^2,&n\equiv1\pmod4,\\
 s^2,&n\equiv3\pmod4
 \end{cases}
 \qquad(s\in\Z\setminus\{0\}).
\]
Consequently, for two primitive integral representatives with
$a,a'>0$, the class in \eqref{eq:sym-projective-T} equals
$[aa'+nbb']$.
\end{proposition}

\begin{proof}
If $a+e=0$, the first quadric gives $c=0$ and the second gives
$nb^2=e^2$.  Since $n>1$ is squarefree, this would be the zero
vector.  Multiplication by $2/(a+e)$ gives the representative
\[
 (a,b,c,e)=(1+h^2,\beta,2h,1-h^2),
\]
because $(a+e)(a-e)=c^2$.  The second quadric becomes
$n\beta^2=\Phi(h)$.  For two normalized representatives,
\eqref{eq:sym-projective-T} is
$(1+h^2)(1+k^2)+n\beta\beta'$, as in
Theorem~\ref{thm:sym-return}.  Its numerator and denominator both
have bidegree $(1,1)$, proving invariance.

If a prime divides $a,c,e$, its square divides $nb^2$ and hence
the prime divides $b$, contradicting primitivity.  Thus
$(a,c,e)$ is a primitive Pythagorean triple: $a$ is odd and the
two legs have opposite parity.  In particular $e^2-c^2$ is odd,
which excludes even $n$ and makes $b$ odd.  When $c$ is even,
the coprime nonnegative integers $(a+e)/2,(a-e)/2$ have square
product $(c/2)^2$, so both are squares; here $n\equiv1\pmod4$.
When $c$ is odd, the coprime positive odd integers $a+e,a-e$
have square product $c^2$, and $n\equiv3\pmod4$.
The product $(a+e)(a'+e')$ is therefore a square.  Finally
$a^2-nb^2=2c^2>0$ at these rational points, so
$aa'>n|bb'|$ and $aa'+nbb'>0$.
\end{proof}

This normalization matters when choosing representatives.  For example,
on $D_{161}$ take
\[
 p=(17,1,8,15),\qquad q=(689,31,400,561).
\]
Then $aa'+161bb'=16704=29\cdot24^2$, whereas replacing $p$ by
$p/16704$ makes this raw value equal to $1$.  In both cases
\eqref{eq:sym-projective-T} is $1044/625=29(6/25)^2$ and
the return is $(49/348,161/696)$.  Its four distances are
$\sqrt{29}(35,115,101,149)/696$.

\subsection{Two pencils of genus-three curves}

The quadric $r_0^2+r_3^2=r_1^2+r_2^2$ also has the ruling
coordinates
\[
 (r_0:r_1:r_2:r_3)
 =(\lambda+\mu:\mu-\lambda:\lambda\mu+1:\lambda\mu-1).
\]
For positive distances with $r_2\ne r_3$, one may take
$\lambda=(r_2+r_3)/(r_0+r_1)$ and
$\mu=(r_0+r_1)/(r_2-r_3)$.  If $s=2/(r_2-r_3)$, the unit-side
condition becomes
\begin{equation}
 \Gamma_\lambda:\quad
 \bigl(2s^2-(1+\lambda^2)(1+\mu^2)\bigr)^2
 =-\Phi(\lambda)\Phi(\mu).
 \label{eq:sym-plane-quartic}
\end{equation}
For a rational point of this curve with $s\ne0$, the return is
\[
 x=\frac12+\frac{2\lambda\mu}{s^2},\qquad
 y=\frac12-\frac{(\lambda^2-1)(\mu^2-1)}{2s^2},
\]
and the four squared distances are the squares of the four displayed
bilinear coordinates divided by $s^2$.  Thus the same strict chamber
selects the original configurations on each fibre.

\begin{theorem}[Quartic pencils and elliptic factors]
\label{thm:sym-quartic-pencil}
For $\lambda\in\Q\setminus\{0,1,-1\}$, the projective
completion of \eqref{eq:sym-plane-quartic} is a smooth plane
quartic.  Put $c_\lambda=-\Phi(\lambda)$ and
\[
 E_{c_\lambda}:\quad y^2=x^3-c_\lambda^2x,
 \qquad
 F_\lambda:\quad v^2=c_\lambda(s^2-4)(s^2-4\lambda^2),
\]
where $(s,v)=(2,0)$ is chosen as origin on $F_\lambda$.
Then there is a $\Q$-isogeny
\[
 \Jac(\Gamma_\lambda)\sim_\Q
 E_{c_\lambda}\times F_\lambda^2.
\]
The other ruling gives a second pencil with the same description.
If the four distances lie in $\sqrt d\,\Q$, the two
$F_\lambda$ factors are replaced by their quadratic twists by $d$.
\end{theorem}

\begin{proof}
Set $A=1+\lambda^2$, $D=(\lambda^2-1)^2$ and
$B=\lambda^4-10\lambda^2+1$.  In the squared projective
coordinates $M=\mu^2$, $S=s^2$, $Z=z^2$, the equation is the conic
\[
 2S^2-2AS(M+Z)+D(M^2+Z^2)-2BMZ=0.
\]
Its symmetric matrix has determinant $-4\Phi(\lambda)^2$.
The discriminants on $M=0$ and $Z=0$ are
$-4\Phi(\lambda)$; on $S=0$ the discriminant is
$-64\lambda^2\Phi(\lambda)$.  They are nonzero, and the
conic contains none of the three coordinate vertices.  Pulling
back by $(\mu:s:z)\mapsto(\mu^2:s^2:z^2)$ therefore gives
a smooth degree-four plane curve.  Equivalently, its degree-four
cover of the conic has six distinct branch points, each of inertia
two, so $2g-2=-8+12=4$.

The involutions changing the signs of $\mu$ and $s$ generate
a Klein four-group.  The quotient by the sign of $s$ is
\[
 C_\lambda:\quad W^2=c_\lambda\Phi(\mu).
\]
It maps nonconstantly to
$E':Y^2=X(X^2-6c_\lambda X+c_\lambda^2)$ by
$(X,Y)=(c_\lambda\mu^2,c_\lambda\mu W)$.
The usual degree-two map
\[
 X_1=\frac{X^2-6c_\lambda X+c_\lambda^2}{X},\qquad
 Y_1=Y\left(1-\frac{c_\lambda^2}{X^2}\right)
\]
has image $Y_1^2=X_1(X_1+4c_\lambda)(X_1+8c_\lambda)$.
Setting $X_1=4(x-c_\lambda)$, $Y_1=8y$ gives
$E_{c_\lambda}$; hence $\Jac(C_\lambda)$ is isogenous to it.

The quotient by the sign of $\mu$ is $F_\lambda$.  Indeed,
with
$v=D\mu^2-B-As^2$, its equation is
$v^2=c_\lambda(s^2-4)(s^2-4\lambda^2)$.
Interchanging the projective coordinates $\mu,z$ preserves the
quartic and conjugates this involution to the third nontrivial
element of the Klein group.  The third quotient is consequently
isomorphic to $F_\lambda$ over $\Q$.
The pullbacks of the three genus-one differentials lie in the three
different nontrivial character spaces.  They span the three-dimensional
space of regular differentials, so the induced homomorphism of
Jacobians is an isogeny.

For a common twist $d$, write the side parameter as $s^2=d t^2$.
The second quotient becomes
$v^2=c_\lambda(d t^2-4)(d t^2-4\lambda^2)$.
Over $\Q(\sqrt d)$ it is identified with $F_\lambda$ by
$s=\sqrt d\,t$.  The nontrivial Galois action changes $s$ to
$-s$, which acts by negation on the Jacobian, as is also seen on
the differential $ds/v$.  This is the quadratic twist by $d$.
\end{proof}
